% ECONOMICS_PROBLEM: {"id": "G14-582", "title": "Limits to arbitrage and anomaly persistence", "jel_code": "G14", "source_id": "OP-0093"}
% !TeX program = xelatex
\documentclass[11pt,a4paper]{article}
\usepackage[margin=23mm]{geometry}
\usepackage{fontspec}
\setmainfont{Latin Modern Roman}
\usepackage{amsmath,amssymb,mathtools}
\usepackage{xcolor,enumitem,booktabs,longtable,array,xurl,hyperref}
\definecolor{muted}{HTML}{53636C}
\hypersetup{colorlinks=true,urlcolor=blue,linkcolor=blue}
\setlength{\parindent}{0pt}
\setlength{\parskip}{5pt}
\newcommand{\R}{\mathbb R}
\newcommand{\E}{\mathbb E}
\newcommand{\N}{\mathbb N}
\newcommand{\ind}{\mathbf 1}
\newcommand{\Var}{\operatorname{Var}}
\newcommand{\Cov}{\operatorname{Cov}}
\newcommand{\supp}{\operatorname{supp}}
\DeclareMathOperator*{\argmin}{arg\,min}
\DeclareMathOperator*{\argmax}{arg\,max}
\newcommand{\entrymeta}[3]{{\sffamily\small\color{muted}\textbf{\texttt{#1}}\quad|\quad #2\par #3\par}}
\newcommand{\Problemref}[1]{\texttt{#1}}
\newcommand{\JELref}[2]{\texttt{#2} (JEL \texttt{#1})}
\begin{document}
\section*{Limits to arbitrage and anomaly persistence}
\textbf{Problem ID:} \texttt{G14-582}\quad \textbf{JEL:} \texttt{G14}\par
% BEGIN ECONOMICS STATEMENT
\label{problem:OP-0093}
\label{atlas:prob:F3}

\noindent\textbf{Atlas ID:} F3; \textbf{original number:} 93; \textbf{field:} Finance. \textbf{Classification:} Editorial research formulation (theoretical/empirical); not a certified unsolved theorem.

\subsection*{Definitions and assumptions}
Fix a tradable asset universe and information filtration. A predictable strategy $w_t$ satisfies specified leverage, short-sale and turnover constraints. Its frictionless, zero-initial-value self-financing gain is $\Pi^{\mathrm{gross}}_{t+1}(w)$. Define its implementable net profit by $\Pi^{\mathrm{net}}_{t+1}(w)=\Pi^{\mathrm{gross}}_{t+1}(w)-c_{t+1}(w)$, where $c_{t+1}(w)\ge0$ records additional execution costs and borrowing spreads relative to the gross strategy's financing benchmark, counted once. Model $\theta$ specifies risk, preferences, constraints and equilibrium selection, and implies a candidate positive pricing kernel $m_{t+1,\theta}$. Assume integrable gains, costs and kernel products. In a frictionless unconstrained pricing model the gross gain satisfies $\mathbb E[m_{t+1,\theta}\Pi^{\mathrm{gross}}_{t+1}(w)]=0$, so the discounted net gain equals minus the discounted costs.

\subsection*{Formal research question}
For a preregistered strategy set $\mathcal W$, independently disciplined model class $\Theta$ and observed population law $P_O$, investigate
\[
 \mathcal A=\{(\theta,w):\theta\in\Theta,\ w\in\mathcal W,
  \mathbb E_{P_O}[m_{t+1,\theta}\Pi^{\mathrm{gross}}_{t+1}(w)]
   \ne\kappa_\theta(w)\},
\]
where $\kappa_\theta(w)=0$ for the frictionless benchmark; a constrained model must derive its nonzero wedge from its economic primitives. The equivalent net-profit restriction subtracts $\mathbb E_{P_O}[m_{t+1,\theta}c_{t+1}(w)]$ from the predicted gross moment. Determine which model violations survive sampling error, multiple testing and prospective evaluation, and whether the associated strategy remains profitable after implementable costs. Latent kernel components require explicit identification assumptions.

\subsection*{Interpretation and scope}
Known trading costs are not mispricing: testing the zero-moment restriction on net profits would reject a correctly priced gross strategy mechanically. A positive mean excess return need not be an arbitrage, which requires nonnegative feasible payoffs almost surely and strictly positive payoffs with positive probability at nonpositive initial cost. Factor alpha is a joint test of the risk benchmark and prices. Allowing an arbitrary pricing kernel or wedge destroys discrimination, so independent restrictions are essential. Publication and strategy selection can cause mechanical decay; specify the discovery process and prospective horizon. Include delistings and realistic execution, and distinguish changes in funding availability from contemporaneous risk or asset-supply shocks.

\subsection*{Source audit and status}
[S1] establishes a funding-constraint mechanism, [S2] surveys limits of arbitrage, and [S3] studies post-publication predictability. These do not classify every anomaly. The editorial agenda concerns independently disciplined risk models, implementable net returns and prospective evidence; established benchmark mechanisms are not recast as open theorems.

\subsection*{References}

\begin{enumerate}[label={[S\arabic*]},leftmargin=*]

\item Andrei Shleifer and Robert W. Vishny (1997). \emph{The Limits of Arbitrage}. Journal of Finance 52(1), 35–55. \url{https://doi.org/10.1111/j.1540-6261.1997.tb03807.x}.
\textit{Locator and role:} Publisher or author abstract / bibliographic record; Funding and delegated-arbitrage constraints. Provides background or a benchmark result; does not establish that the editorial question remains unresolved in 2026.

\item Denis Gromb and Dimitri Vayanos (2010). \emph{Limits of Arbitrage}. Annual Review of Financial Economics 2, 251–275. \url{https://doi.org/10.1146/annurev-financial-073009-104107}.
\textit{Locator and role:} Publisher or author abstract / bibliographic record; Published theoretical survey; working-paper title adds The State of the Theory. Provides background or a benchmark result; does not establish that the editorial question remains unresolved in 2026.

\item R. David McLean and Jeffrey Pontiff (2016). \emph{Does Academic Research Destroy Stock Return Predictability?}. Journal of Finance 71(1), 5–32. \url{https://doi.org/10.1111/jofi.12365}.
\textit{Locator and role:} Publisher or author abstract / bibliographic record; Post-discovery and post-publication changes in return predictability. Provides background or a benchmark result; does not establish that the editorial question remains unresolved in 2026.

\end{enumerate}

\subsection*{Common conventions: Source mathematical conventions}

Unless an entry states otherwise, agent and firm sets are finite. State and action spaces are measurable, strategies and outcome maps are measurable, probability laws are specified, and expectations exist for the targets under discussion. Infinite-horizon discount factors lie in $(0,1)$ and discounted objectives are finite. Norms, tolerances and complexity input sizes must be specified for computational comparisons. Constants or primitives that appear locally are inputs, not empirical facts asserted by this catalogue.

\textbf{Identification.} A statistical model $m\in\mathcal M$ implies an observable law $P_m$ and target $\theta(m)$. At law $P$, its identified set is
\[
\Theta(P;\mathcal M)=\{\theta(m):m\in\mathcal M,\ P_m=P\}.
\]
Point identification holds when this set is a singleton, or equivalently when $P_{m_1}=P_{m_2}$ implies $\theta(m_1)=\theta(m_2)$ for all compatible models. Sharp bounds mean bounds attained, or approached when the boundary is not attained, by compatible models. Writing this set is a definition, not a solution: the substantive task is to establish informative restrictions and derive or estimate the set. An unrestricted-confounding model may give only trivial bounds.

\textbf{Inference.} Identification, estimability and valid uncertainty quantification are separate questions. Fix $\alpha\in(0,1)$. A confidence procedure $C_n$ over a declared law class $\mathcal P$ has asymptotic uniform coverage at level $1-\alpha$ when
\[
\liminf_{n\to\infty}\inf_{P\in\mathcal P}\Pr_P\{\theta(P)\in C_n\}\geq1-\alpha.
\]
For set-identified targets the coverage event must be defined separately, for example coverage of the full identified set. If an entry uses identification-indexed models instead of uniquely indexed laws, the same coverage convention is applied over those models. Pointwise limits do not establish uniform validity, and asymptotic validity does not guarantee finite-sample precision. Sample sizes, dependence structures and weak-identification sequences are part of each application.

\textbf{Counterfactuals.} $Y_i(a)$ is a potential outcome under a specified intervention, including its timing, implementation and versions. It is not $Y_i$ conditional on voluntarily choosing $a$. A full intervention vector permits interference; shorthand scalar treatment notation does not silently impose no interference. Assignment, selection, attrition, anticipation and measurement must be specified. A claim about an instrument's local effect is not automatically a population policy effect.

\textbf{Economic equilibrium.} A counterfactual model includes best responses, constraints, feasibility, market clearing, state transitions and expectations consistent with the stated information. When several equilibria exist, either a declared selector is an input or the target is set-valued. Comparisons across policy regimes must use the stated selection convention. Reduced-form relationships are not presumed invariant after an equilibrium-changing intervention.

\textbf{Optimization and welfare.} Optimization is over the stated feasible set. If attainment is not established, $\sup$ or $\inf$ and $\varepsilon$-optimal choices replace an asserted maximizer or minimizer. For example, with value $V=\sup_{a\in\mathcal A}W(a)<\infty$, an $\varepsilon$-optimal set is $\{a\in\mathcal A:W(a)\geq V-\varepsilon\}$ for $\varepsilon>0$. Benefits, costs, risk and health or environmental outcomes can be aggregated only using declared units and conversion weights. Interpersonal utility weights, the treatment of fiscal transfers and foreign profits, discounting, and the behavioral welfare criterion are explicit inputs. A comparative-static derivative additionally requires existence and appropriate differentiability of the selected counterfactual.

\textbf{Sources.} Local labels [S1], [S2], and so forth restart in every question. Locators identify the relevant abstract, model, section or result at the level actually checked; a bibliographic or abstract-level verification is not represented as inspection of an entire proof. Working papers are distinguished from later publications and changing versions. Source summaries are paraphrased. All URL checks and editorial formulations refer to the audit date stated above.
% END ECONOMICS STATEMENT
\end{document}
